\section{Cost-shape generalisation: code generation on HumanEval}
\label{app:humaneval}

The cost laws of \S\ref{sec:laws} are derived from structured log
extraction (BGL). Every experiment in the main text --- chains,
bounded deadline, wide schema, perturbation --- runs on the same task.
This appendix tests whether the line-vs-staircase cost shape
generalises to a second task: code generation on
HumanEval~\cite{humaneval} (164 problems, Python function completion).

Code generation has the properties that favour block-diffusion: output
is long (50--300 tokens), truncation is catastrophic
(\texttt{SyntaxError} = pass@1 = 0), schema validity is enforced by
Python syntax, and batch-1 is the natural mode (IDE autocomplete,
runtime codegen).

\subsection{Baseline: cost law holds on code}

Both arms were run serial (batch-1) on the pinned stack
(\S\ref{sec:threats}), $n{=}164$ problems.

\begin{table}[h]
\centering
\caption{HumanEval baseline (serial, batch-1, $n{=}164$). The cost-shape
law holds on code: AR scales linearly at 6.4\,ms/token (vs 7.2 on logs),
DG is a flat staircase. AR is 2.9$\times$ slower at natural output
(vs 2.5$\times$ on logs) --- the ratio is in the same range.}
\label{tab:humaneval}
\begin{tabular}{@{}lrrrrr@{}}
\toprule
Arm & pass@1 & code valid\% & p50 (ms) & mean tokens & ms/token \\
\midrule
DG K=8  & 0.287 & 31.1 & 280 & 99.5  & --- (staircase) \\
AR-30B  & 0.720 & 93.9 & 800 & 124.1 & 6.4 \\
AR-3B   & 0.573 & 99.4 & 560 & 107.7 & 5.2 \\
\bottomrule
\end{tabular}
\end{table}

AR-30B's per-token slope on code is close to its log-extraction slope.
DG's p50 is a single-canvas staircase cost, independent of output
length. The AR/DG ratio matches the log-extraction ratio --- the cost
shape generalises.

DG's lower pass@1 is expected: DiffusionGemma was trained as a
log-extraction model, not a code model. What matters is not absolute
quality but whether the \emph{shape} (line vs staircase) transfers.

\subsection{Deadline-bounded pass@1: AR collapses, DG stays flat}

Under a tightening deadline, AR's \texttt{max\_tokens} is set to
$\lfloor D / \text{slope} \rfloor$ (what it can generate in time).
DG always generates a full canvas regardless of $D$.

\begin{table}[h]
\centering
\caption{pass@1 under tightening deadline ($n{=}164$, serial). AR
collapses to 0.085 at $D{=}200$\,ms (truncated code = SyntaxError).
DG stays at 0.287 regardless of $D$ (always emits a complete canvas).
Crossover at $D \leq 400$\,ms: DG wins despite being worse at code.}
\label{tab:deadline}
\begin{tabular}{@{}rrrrrr@{}}
\toprule
$D$ (ms) & AR pass@1 & AR valid\% & AR p50 (ms) & DG pass@1 & DG p50 (ms) \\
\midrule
200  & 0.085 & 22.6 & 230 & 0.287 & 279 \\
400  & 0.189 & 33.5 & 434 & 0.287 & 280 \\
600  & 0.274 & 45.7 & 637 & 0.287 & 280 \\
800  & 0.311 & 54.3 & 801 & 0.287 & 280 \\
$\infty$ & 0.720 & 93.9 & 800 & 0.287 & 280 \\
\bottomrule
\end{tabular}
\end{table}

AR collapses because a truncated function is a \texttt{SyntaxError},
not just a lower-quality answer. On logs, truncation means missing
JSON keys (bad but recoverable); on code, it means a crash. This makes
the graceful-degradation property (drop $K$, keep completeness)
visibly more valuable on code than on logs.

\subsection{Multi-turn correction chain: per-hop latency stays flat}

A code correction chain tests whether DG's per-hop latency stays
constant across \emph{variable-content} rounds (generate $\to$ execute
$\to$ fix $\to$ re-execute), not just repeated identical calls. Up to
5 rounds per problem.

\begin{table}[h]
\centering
\caption{Multi-turn correction chain ($n{=}164$, serial, max 5 rounds).
DG's per-round latency stays flat ($+6\%$); AR's grows ($+21\%$) as
fix prompts accumulate context. Both saturate by round 3.}
\label{tab:chain}
\begin{tabular}{@{}lrrrr@{}}
\toprule
Arm & R1 pass@1 & cumulative & R1 p50 (ms) & total p50 (ms) \\
\midrule
DG K=8  & 0.854 & 0.957 & 287 & 305 \\
AR-30B  & 0.762 & 0.896 & 511 & 620 \\
\bottomrule
\end{tabular}
\end{table}

DG's total p50 (305\,ms) is barely above its single-round p50
(287\,ms), because most problems solve in round 1 and the few that
need a fix round add little. AR's total (620\,ms) is $1.2\times$ its
single-round latency (511\,ms), because the fix round re-prefills a
longer accumulated prompt. The staircase shape holds across
variable-content hops, confirming the chain latency claim from
\S\ref{sec:survives} with a real multi-turn workflow.
